\documentclass[a4paper,12pt]{article}
\usepackage[T2A]{fontenc}
\usepackage[utf8]{inputenc}
\usepackage[russian]{babel}
\usepackage{amsmath,amssymb,geometry,booktabs,tabularx,xcolor,tikz,fancyhdr,hyperref}
\geometry{margin=23mm,headheight=15pt}
\definecolor{accent}{HTML}{008080}
\hypersetup{colorlinks=true,linkcolor=accent,pdftitle={Ревью домашней работы}}
\setlength{\parindent}{0pt}
\setlength{\parskip}{5pt}
\pagestyle{fancy}\fancyhf{}\fancyhead[L]{Ревью домашней работы}\fancyhead[R]{Анастасия Павлова}\fancyfoot[C]{\thepage}
\newcommand{\parthead}[1]{\par\smallskip\textbf{#1}\par\nopagebreak[4]}
\newcommand{\err}[1]{\par\textcolor{red!75!black}{\textbf{Ошибка:}} #1\par}
\newcommand{\idea}[1]{\par\textcolor{blue!60!black}{\textbf{Ключевая идея:}} #1\par}
\newcommand{\result}[1]{\par\textcolor{accent}{\textbf{Итог:}} #1\par}
\newcommand{\task}[2]{\clearpage\phantomsection\label{task:#1}\section*{\textcolor{accent}{Задание №#1. #2}}}
\begin{document}
\begin{titlepage}\centering
\vspace*{25mm}
{\Large\textcolor{accent}{МАТЕМАТИКА}}\par\vspace{15mm}
{\Huge Ревью домашней работы}\par\vspace{12mm}
{\Large\textcolor{accent}{Линейные функции, векторы,\\ совместная работа и растворы}}\par
\vfill
{\large Анастасия Павлова}\par\vspace{8mm}
25 ноября 2026 г.\par\vfill
{\small Лёвин Артём Александрович\par эксклюзивно для Анастасии Павловой}\par\vspace{14mm}
\begin{tikzpicture}\draw[accent,thick] (0,0) .. controls (2,0.7) and (3,-0.7) .. (5.5,0) .. controls (8,0.7) and (9,-0.7) .. (11,0);\end{tikzpicture}
\end{titlepage}
\section*{Сводная таблица замечаний}
Нажмите на номер задания для перехода к разбору. В таблице отмечены вычислительные ошибки, ошибки записи и неполнота ответа.
\par\smallskip\noindent%
{\small\renewcommand{\arraystretch}{1.5}
\begin{tabularx}{\dimexpr\linewidth-4pt\relax}{@{}l >{\raggedright\arraybackslash}p{26mm} >{\raggedright\arraybackslash}X >{\raggedright\arraybackslash}p{29mm} >{\raggedright\arraybackslash}p{32mm}@{}}
\toprule № & Тема & Суть замечания & Ваш ответ & Правильный ответ\\\midrule
\hyperref[task:3]{3} & Линейная функция & Формула функции не записана отдельно & $k=-2$, $b=3$, $9$ & $f(x)=-2x+3$, $f(-3)=9$\\
\hyperref[task:4]{4} & Пересечение прямых & Ошибка при решении линейного уравнения & $(-3;9)$ & $(1;5)$\\
\hyperref[task:6]{6} & Векторы & Косинус записан с единицей угла & $0^\circ$; $90^\circ$ & $0$; $90^\circ$\\
\hyperref[task:7]{7} & Совместная работа & Пропущен знак коэффициента при подстановке в дискриминант & 6 ч; 12 ч & 6 ч; 12 ч\\
\hyperref[task:10]{10} & Растворы & Ошибка сложения масс соли & $22\%$ & $32\%$\\\bottomrule
\end{tabularx}}
\clearpage\thispagestyle{plain}\null\vfill\begin{center}{\Huge\textcolor{accent}{Разбор ошибок}}\end{center}\vfill\null
\task{3}{Линейная функция}
\parthead{Текст задания}
График $f(x)=kx+b$ проходит через $A(-2;7)$ и $B(2;-1)$. Найдите формулу $f(x)$ и значение $f(-3)$.
\parthead{Ваше решение и ответ}
Вы составили систему, получили $b=3$, $k=-2$, затем вычислили $f(-3)=-2\cdot(-3)+3=9$. Оба коэффициента и числовое значение найдены верно.
\err{Неполное оформление ответа: отсутствует отдельная запись формулы функции.}
\parthead{Где возникает пропуск и почему он важен}
Последний корректный шаг --- определение коэффициентов $k=-2$ и $b=3$. Далее можно вычислить значение в заданной точке, однако условие также требует явно записать зависимость для произвольного $x$. Число $9$ характеризует только одно значение аргумента.
\idea{Формула линейной функции получается подстановкой найденных коэффициентов в общий вид $f(x)=kx+b$.}
\parthead{Исправление и полное решение}
Подстановка координат двух точек даёт систему:
\[\begin{cases}-2k+b=7,\\2k+b=-1.\end{cases}\]
Сложим уравнения: $2b=6$, откуда $b=3$. Подставим в первое: $-2k+3=7$, $-2k=4$, $k=-2$. Поэтому
\[f(x)=-2x+3,\qquad f(-3)=6+3=9.\]
\parthead{Проверка}
$f(-2)=4+3=7$, $f(2)=-4+3=-1$. Обе исходные точки принадлежат графику.
\result{$f(x)=-2x+3$; $f(-3)=9$. Вычислительную часть решения можно сохранить полностью.}
\parthead{Задача на закрепление}
График $f(x)=kx+b$ проходит через $A(-1;5)$ и $B(3;-3)$. Найдите формулу функции и значение $f(-4)$.
\task{4}{Пересечение прямых}
\parthead{Текст задания}
Прямая $y=kx+b$ проходит через $(-1;1)$ и $(3;9)$. Найдите точку её пересечения с прямой $y=-x+6$.
\parthead{Ваше решение и ответ}
Вы нашли $k=2$, $b=3$ и правильно записали $-x+6=2x+3$. Затем получили $x=-3$, $y=9$. Ваш ответ: $(-3;9)$.
\err{Неверное решение линейного уравнения при поиске абсциссы пересечения.}
\parthead{Где именно возникает ошибка}
Последняя верная строка: $-x+6=2x+3$. Следующая строка $x=-3$ неверна. Промежуточные преобразования отсутствуют, поэтому конкретную причину ошибочного вычисления установить нельзя.
\parthead{Почему это неверно}
При переносе слагаемого в другую часть уравнения меняется его знак. Это следствие прибавления или вычитания одной и той же величины из обеих частей. После приведения подобных необходимо разделить обе части на коэффициент при $x$.
\idea{В точке пересечения обе формулы дают одно и то же значение $y$ при одном и том же $x$.}
\parthead{Исправление и полное решение}
Из условий $-k+b=1$, $3k+b=9$ вычтем первое уравнение из второго: $4k=8$, $k=2$, $b=3$. Искомая прямая имеет вид $y=2x+3$. Приравняем правые части:
\begin{align*}-x+6&=2x+3,\\6-3&=2x+x,\\3&=3x,\\x&=1.\end{align*}
Тогда $y=2\cdot1+3=5$.
\parthead{Проверка}
При $x=1$ обе прямые дают $5$: $-1+6=5$ и $2+3=5$. При $x=-3$ первая даёт $-3$, вторая $9$ (по формулам $2x+3$ и $-x+6$ соответственно), поэтому прежняя точка не принадлежит обеим прямым.
\result{$(1;5)$.}
\parthead{Задача на закрепление}
Прямая проходит через $(-2;0)$ и $(2;8)$. Найдите точку её пересечения с прямой $y=-x+7$.
\task{6}{Косинус угла между векторами}
\parthead{Текст задания}
Даны $A(0;1)$, $B(2;2)$, $C(-1;3)$. Найдите $\cos\angle BAC$ и определите величину угла.
\parthead{Ваше решение и ответ}
Вы правильно нашли $\overrightarrow{AB}=(2;1)$, $\overrightarrow{AC}=(-1;2)$ и нулевое скалярное произведение. В итоговой записи косинуса поставлен знак градуса: $0^\circ$. Величина угла записана верно: $90^\circ$.
\err{Косинус записан как величина угла.}
\parthead{Где именно возникает ошибка}
Вычисление приводит к числу $0$. Ошибка появляется в итоговой записи $\cos\angle BAC=0^\circ$. Знак градуса относится к углу. Косинус является безразмерным числом из отрезка $[-1;1]$.
\idea{Различайте значение косинуса и величину угла. Нулевой косинус соответствует прямому углу.}
\parthead{Исправление и полное решение}
Вычтем координаты начала каждого вектора из координат его конца:
\[\overrightarrow{AB}=(2-0;2-1)=(2;1),\qquad
\overrightarrow{AC}=(-1-0;3-1)=(-1;2).\]
Скалярное произведение равно $2\cdot(-1)+1\cdot2=0$. Длины обоих векторов равны $\sqrt5$, поэтому произведение длин равно $5$. По формуле
\[\cos\angle BAC=\frac{\overrightarrow{AB}\cdot\overrightarrow{AC}}{|\overrightarrow{AB}|\,|\overrightarrow{AC}|}=\frac{0}{5}=0.\]
На отрезке от $0^\circ$ до $180^\circ$ косинус обращается в ноль при $90^\circ$.
\parthead{Проверка}
Векторы ненулевые, их скалярное произведение равно нулю; следовательно, они перпендикулярны.
\result{$\cos\angle BAC=0$, $\angle BAC=90^\circ$.}
\parthead{Задача на закрепление}
Даны $A(1;1)$, $B(4;2)$, $C(0;4)$. Найдите косинус угла $BAC$ и величину угла. Запишите оба результата с корректными обозначениями.
\task{7}{Совместная работа}
\parthead{Текст задания}
Первый мастер может выполнить работу на 6 часов быстрее второго. Вместе они выполняют её за 4 часа. Сколько часов потребуется каждому по отдельности?
\parthead{Ваше решение и ответ}
Вы составили уравнение $\frac1x+\frac1{x+6}=\frac14$ и правильно получили $x^2-2x-24=0$. В строке дискриминанта записано $(-2)^2-4\cdot1\cdot24=4+96$. Далее найдено $x=6$. Ответ: 6 и 12 часов.
\err{В записи дискриминанта потерян отрицательный знак свободного члена. Численный ответ при этом верен.}
\parthead{Где именно возникает ошибка и почему}
Для уравнения $ax^2+bx+c=0$ формула $D=b^2-4ac$ требует подстановки коэффициентов вместе с их знаками. Здесь $c=-24$. Если подставить $24$, выражение даст $4-96=-92$. Положительный результат $100$ получается при корректной подстановке отрицательного числа в скобках.
\idea{Сначала явно выпишите $a$, $b$, $c$ со знаками; затем подставьте их в формулу.}
\parthead{Исправление и полное решение}
Пусть $x>0$ --- время первого мастера в часах; второй работает $x+6$ часов. Их производительности складываются:
\[\frac1x+\frac1{x+6}=\frac14.\]
Умножим на $4x(x+6)$, который при $x>0$ отличен от нуля:
\[4(x+6)+4x=x(x+6)\quad\Longrightarrow\quad x^2-2x-24=0.\]
Коэффициенты $a=1$, $b=-2$, $c=-24$. Тогда
\[D=(-2)^2-4\cdot1\cdot(-24)=100,\qquad x_{1,2}=\frac{2\pm10}{2}.\]
Получаем $x_1=6$ и $x_2=-4$. Второй корень исключаем: время положительно. Второму мастеру требуется $6+6=12$ часов.
\parthead{Проверка}
$\frac16+\frac1{12}=\frac14$, а $12-6=6$ часов. Оба условия выполнены.
\result{Первый мастер --- 6 ч, второй --- 12 ч.}
\parthead{Задача на закрепление}
Первый мастер выполняет заказ на 8 часов быстрее второго. Вместе они выполняют заказ за 3 часа. Найдите время каждого по отдельности.
\task{10}{Концентрация раствора}
\parthead{Текст задания}
Из 6 кг 20\%-го раствора выпарили 1 кг воды, затем добавили 5 кг 40\%-го раствора. Найдите концентрацию соли в полученной смеси.
\parthead{Ваше решение и ответ}
Вы верно составили таблицу: исходная масса соли $1{,}2$ кг, в испарившейся воде соли нет, добавленная масса соли $2$ кг, итоговая масса раствора $10$ кг. Получено уравнение, эквивалентное $1{,}2+2=10x$, но в следующей строке записано $2{,}2=10x$. Ваш ответ: $22\%$.
\err{Арифметическая ошибка: $1{,}2+2=3{,}2$, а не $2{,}2$.}
\parthead{Где именно возникает ошибка и почему}
До составления уравнения решение правильное. При сложении десятичной дроби и целого числа выравниваем разряды: $1{,}2+2{,}0=3{,}2$. Неверная масса соли уменьшила результат на 10 процентных пунктов.
\idea{Отдельно отслеживайте общую массу раствора и массу растворённой соли. При выпаривании чистой воды масса соли сохраняется.}
\parthead{Исправление и полное решение}
Масса соли в исходном растворе: $6\cdot0{,}20=1{,}2$ кг. После выпаривания остаётся $6-1=5$ кг раствора с той же массой соли. В добавленном растворе содержится $5\cdot0{,}40=2$ кг соли.
\begin{center}\small\renewcommand{\arraystretch}{1.25}
\begin{tabularx}{\linewidth}{@{}Xrr@{}}\toprule Стадия & Раствор, кг & Соль, кг\\\midrule Исходный раствор & 6 & 1,2\\ Выпаривание воды & $-1$ & 0\\ Добавление раствора & 5 & 2\\\midrule Итог & 10 & 3,2\\\bottomrule\end{tabularx}\end{center}
Итак, $10x=3{,}2$, $x=0{,}32$. В процентах это $32\%$.
\parthead{Проверка}
После выпаривания концентрация равна $1{,}2/5=24\%$. Смешиваются равные массы растворов с концентрациями $24\%$ и $40\%$; их среднее равно $32\%$. Это независимо подтверждает результат.
\result{$32\%$. Табличная модель составлена правильно.}
\parthead{Задача на закрепление}
Из 10 кг 30\%-го раствора выпарили 2 кг воды, затем добавили 4 кг 15\%-го раствора. Найдите концентрацию соли в полученной смеси.
\vfill\begin{center}\small Лёвин Артём Александрович\par эксклюзивно для Анастасии Павловой\end{center}
\end{document}
